\documentclass{situationist-preprint}

\title{What \enquote{Barely Worth Living} Means}
\msid{ETH-09-b}
\msfamily{ETH-09}
\msdiscipline{Normative and Applied Ethics}
\mstype{Companion manuscript}
\mskeywords{repugnant conclusion; lives barely worth living; neutral level; worth starting; worth continuing}

\begin{document}
\maketitle

\begin{abstract}
ETH-09-a argues that the repugnant conclusion is true and not repugnant, because the intuition against it misdescribes lives barely worth living. This manuscript examines the notion itself. We distinguish three ambiguities in \enquote{barely worth living}: between the value of a whole life and the quality of its moments, between a life worth starting and a life worth continuing, and between small value in absolute terms and small value relative to the lives compared. Resolved in the way the population axiology requires, the phrase picks out whole lives, assessed as worth starting, whose value is small relative to the lives in the better-off population. So understood, the class includes many lives like those actually lived, and the residual case of the drabbest such lives does not restore the intuition.
\end{abstract}

\section{Three ambiguities}

ETH-09-a defines the value of a life as the value that idealized reflective endorsement assigns to it, counting each component of the good realized in it (Definition~6.1), and calls a life barely worth living if that value is positive and small relative to the values of the lives with which it is being compared (Definition~6.2). The zero of the scale is the neutral level of Definition~2.1, the level of a life neither worth living nor not. Three ambiguities in ordinary uses of the phrase separate these definitions from the picture that drives the intuition.

\paragraph{Lives and moments.} A life may be barely worth living because each of its moments is barely good, or because its moments are very good and very bad in a balance that is narrowly favorable. \citet{Parfit1986} illustrated the first kind with lives of muzak and potatoes. The definitions concern only the value of the whole life, and are indifferent between the two kinds. The intuition, by contrast, is formed almost entirely on the first: we imagine lives at a low and even level, and not lives of great love and great loss, of achievement and defeat, whose goods narrowly exceed their evils.

\paragraph{Starting and continuing.} A life may be worth continuing without being worth starting. \citet[ch.~2]{Benatar2006} draws the distinction for his own purposes, and it holds independently of them: a person already embarked on a life has interests, attachments and projects that give her reason to go on, and the bar a life must clear for continuing it to be worth while is lower than the bar for bringing it about. In the repugnant conclusion the question is which populations it is better to bring about. The relevant zero is therefore the neutral level for starting a life, and a life barely worth living is one barely worth starting. When we picture lives that their subjects would reluctantly continue, we are picturing lives that may lie at or below the relevant zero, and our sense that a world of them is bad may be correct about them and irrelevant to the conclusion.

\paragraph{Absolute and relative smallness.} Definition~6.2 of ETH-09-a is comparative. In the repugnant conclusion the lives of Z are small in value relative to the lives of A, which are by hypothesis very good. If A is a population of lives far better than any yet lived, a life small relative to them may be as good as a fortunate actual life. Theorem~3.2 of ETH-09-a holds for every positive level $z$, so the conclusion covers such cases as well as cases near the neutral level. The phrase \enquote{barely worth living} suggests only the latter.

\section{The neutral level}

The second ambiguity locates the zero; it does not yet say where it lies. We take it to lie where idealized reflective endorsement, considering the whole life from the standpoint of whether it is better that it be lived, neither favors nor disfavors its being lived. \citet{Broome2004} argues that the neutral level for existence is vague. We accept that it may be: Definition~2.1 of ETH-09-a requires a ratio scale with a zero, and vagueness in the zero yields a range of admissible scales, on each of which Theorems~3.2 and~4.2 of ETH-09-a hold. The conclusion is then true on every admissible precisification, which is all that truth requires.

What the neutral level is not is equally important. It is not the level of a life one would be content to live, nor that of a life one would be willing to trade one's own for. These are the anchored and misplaced zeros that ETH-09-a identifies as sources of the intuition. Nor is it the level below which a life contains more pain than pleasure, since on the plural account of ETH-01 hedonic balance is only one component of a life's value.

\begin{proposition}\label{p:class}
Given the three resolutions above, the lives barely worth living in Z include lives with substantial goods, such as close relationships, partial achievements and real knowledge, whose evils nearly equal them.
\end{proposition}

\begin{proof}
By the first resolution, the class is fixed by the value of whole lives, so any distribution of goods and evils within a life whose overall value is small and positive qualifies. By Definition~6.1 of ETH-09-a, relationships, achievement and knowledge are components of that value; a life that contains them in substantial measure, together with comparable illness, grief and failure, has a small positive value. Such lives are therefore among those the conclusion concerns.
\end{proof}

\citet{Tannsjo2002} argues that most actual lives are of this kind. We need not accept so strong an empirical claim. It is enough that the class includes such lives, for then the intuition formed on the drab picture is not an intuition about the class.

\section{Objection and reply}

\begin{objection}
The repugnant conclusion holds for every positive level, so it holds for Z populations made up entirely of drab lives. The intuition may be directed at those, and the richer lives of Proposition~\ref{p:class} do nothing to defuse it.
\end{objection}

\begin{reply}
The conclusion does cover such populations. But the remaining sources of the intuition identified by ETH-09-a apply to them with full force. A life of muzak and potatoes is pictured from a standpoint far above it; its zero is easily misplaced; and the number of such lives needed to outweigh A is one that intuition cannot represent. Moreover, on the plural account, a life that lacks knowledge, achievement and relationships altogether is of uncertain sign; if it is below the neutral level for starting, it falls outside the conclusion. Once these factors are removed, what remains is a preference for quality over quantity, which ETH-09-a treats as a genuine but outweighed consideration.
\end{reply}

\section{Conclusion}

\enquote{Barely worth living} denotes whole lives, assessed as worth starting, whose value is small relative to the lives in the better population. The class is not a class of drab lives, and it is not a class of lives near misery. Read correctly, the phrase describes much of human life as it has been lived, and a vast population of such lives is not something reflection finds repugnant.

\begin{thebibliography}{9}\small
\bibitem[Benatar(2006)]{Benatar2006} Benatar, D. (2006). \emph{Better Never to Have Been: The Harm of Coming into Existence}. Oxford University Press.
\bibitem[Broome(2004)]{Broome2004} Broome, J. (2004). \emph{Weighing Lives}. Oxford University Press.
\bibitem[Parfit(1986)]{Parfit1986} Parfit, D. (1986). Overpopulation and the quality of life. In P. Singer (ed.), \emph{Applied Ethics}, 145--164. Oxford University Press.
\bibitem[T{\"a}nnsj{\"o}(2002)]{Tannsjo2002} T{\"a}nnsj{\"o}, T. (2002). Why we ought to accept the repugnant conclusion. \emph{Utilitas}, 14(3), 339--359.
\end{thebibliography}

\end{document}
